\documentclass[10pt]{article}

\usepackage[margin=0.78in]{geometry}
\usepackage{amsmath}
\usepackage{amssymb}
\usepackage{booktabs}
\usepackage{enumitem}
\usepackage{longtable}
\usepackage[table]{xcolor}
\usepackage{url}

\newcommand{\packageid}{SCI-MAP-002}
\newcommand{\appsha}{9aae0e669384c5c0c0dda93debc194d6b8dac787}
\newcommand{\corecommit}{bccc2614663c95aee66ed9b21b6025f5fb4e5f04}
\newcommand{\tracecommit}{0d80c4df9d4694634c4c77830c7701b07af577b1}
\newcommand{\jinc}{\operatorname{jinc}}
\newcommand{\corelabel}{\textbf{DERIVATION}}
\newcommand{\srclabel}{\textbf{SOURCE FACT}}
\newcommand{\inflabel}{\textbf{INFERENCE}}
\newcommand{\runlabel}{\textbf{UNVERIFIED RUNTIME}}
\newcommand{\handofflabel}{\textbf{HANDOFF}}

\title{\packageid{} Scientific-Contract Audit\\
JINC Gridding, Normalization, Response, Support, and Product Interface}
\author{Independent scientific-contract audit}
\date{2026-08-03}

\begin{document}
\maketitle

\begin{center}
\begin{tabular}{@{}ll@{}}
\toprule
Governing application identity & \texttt{\appsha} \\
Frozen independent core & \texttt{\corecommit} \\
Independent-core SHA-256 & \texttt{2c1f9ff95f65422a098846f747ed165d5aeddc5bedd854678bfa7faeebba4e24} \\
Frozen static trace & \texttt{\tracecommit} \\
Static-trace SHA-256 & \texttt{37be33eb76dad4d916dd2772cb32f6570fef7e58ebbde0cf2de1add189d73340} \\
Audit branch & \texttt{codex/audit-sci-map-002} \\
Audit mode & read-only source audit; documentation artifacts only \\
\bottomrule
\end{tabular}
\end{center}

\section{Executive disposition}

\textbf{Verdict: amend; implementation nonconformant to the frozen proposed
contract; validation incomplete; production status remains existing-use-only.}
This verdict is bounded to JINC.  It neither changes the established
implementation nor authorizes a repair, numerical parameter campaign, broader
production use, or a conclusion about naive mapmaking.

The governing source contains the correct three raw algebraic accumulators for
the frozen constant-field gridding estimator,
\begin{equation}
  N_p=\sum_s q_s c_{sp}d_s,\qquad
  C_p=\sum_s q_s c_{sp},\qquad
  Q_p=\sum_s q_s c_{sp}^{2},
  \label{eq:raw}
\end{equation}
and the observation-map finalizer publishes
\begin{equation}
  \widehat m_p=N_p/C_p,\qquad
  W_{p,\mathrm{formal}}=C_p^2/Q_p
  \label{eq:final}
\end{equation}
on retained support.  Negative coefficients keep their signs in $N_p$ and
$C_p$ and enter $Q_p$ quadratically.  The source also constructs the raw
kernel numerator $K_p=\sum_s q_sc_{sp}h_s$ and divides it by the same $C_p$;
for the visible analytic RTC templates, $h_s$ is explicitly generated as a
unit-source response.  These are material conforming subcontracts.

They do not close the package.  The source realizes a square, point-sampled
footprint rather than the frozen radial, pixel-response-averaged contract;
uses fixed absolute $|C_p|$ and $Q_p$ numerical scales that are not invariant
under a change of signal unit; records coefficient-squared time as
``coverage''; and can publish \path{coverage_bool_*}=1 for zero or non-finite
weight.  Shape-vector elements are checked for finiteness but not their
physical domain, and the coefficient is not checked for finiteness before
deposition.  Provenance records requested/effective configuration and coarse
completion state but not the realized normalization, support, array
selection, parallel policy, kernel source, cancellation diagnostics, or
product joins required by the frozen core.

The only incoming post-core handoff is explicitly naive-only.  It supplies no
JINC runtime evidence.  Static reachability shows both ordinary and detector-
beammap JINC paths, including concurrent policies, but there is no exact-SHA
runtime evidence for response, product values, or sequential/parallel
acceptance.  A bounded protocol is proposed in Section~\ref{sec:protocol}; it
has not been requested, run, or interpreted.

\section{Authority, chronology, and epistemic separation}

The independent mathematical reference is
\path{doc/audits/packages/SCI-MAP-002_INDEPENDENT_CORE.tex}, frozen and
committed before any JINC implementation or post-core evidence was opened.
Its SHA-256 is stated above.  The phase-two source-to-equation matrix is
\path{doc/audits/packages/SCI-MAP-002_STATIC_SOURCE_TRACE.md}; it was frozen
separately and is preliminary evidence, not the present verdict.

This report uses the following labels strictly:

\begin{itemize}[leftmargin=*]
  \item \corelabel{} is mathematics or a required scientific identity from
        the frozen independent core.
  \item \srclabel{} is directly visible at \texttt{\appsha}; line ranges are
        cited.  The audit branch differs from the governing SHA only by audit
        documentation commits, and the examined application paths have no
        diff from \texttt{\appsha}.
  \item \inflabel{} is a conclusion that follows from combined source facts;
        it is not described as observed runtime behavior.
  \item \runlabel{} names behavior that static reading cannot establish.
  \item \handofflabel{} is bounded incoming evidence and retains the scope
        assigned by its author.
\end{itemize}

No Citlali compilation, test, reduction, numerical probe, Unity access,
network access, delegation, or independent review was performed.  The only
allowed execution for this package is documentation-only LaTeX verification.

\subsection{Post-core handoff disposition}

\handofflabel{} \path{SCI-MAP-002-XAUD-001} states that all seven reductions
in the completed MAP evidence bundle selected \texttt{mapmaking.method naive}.
It explicitly does not establish any JINC implementation, response,
normalization, parameter-selection, edge, parallel-equivalence, uncertainty,
product, or production claim.  This audit acknowledges and incorporates only
the naive-control distinction.  No JINC finding is closed or downgraded by
that handoff, and the ordinary-naive positive-coefficient predicate is not
reused here.

\section{Bounded source expansion}

The phase-two roots were exactly
\path{include/citlali/core/mapmaking/jinc_mm.h},
\path{include/citlali/core/pipeline/mapmaking_dispatch.h}, and
\path{include/citlali/core/pipeline/mapmaking_method_config.h}.  Final audit
inspection followed only direct definitions, callers, publishers, and
consumers needed to answer the frozen questions.  The complete per-file hash
and reason record is in
\path{doc/audits/evidence/SCI-MAP-002_LOCAL_EVIDENCE_2026-08-03.yaml}.
The expansion was grouped as follows:

\begin{longtable}{@{}p{0.21\textwidth}p{0.71\textwidth}@{}}
\toprule
Question & Direct paths opened and reason \\
\midrule
\endhead
Final normalization and masks &
\path{include/citlali/core/mapmaking/map.h},
\path{src/citlali/core/mapmaking/map.cpp}, and
\path{include/citlali/core/utils/utils.h}: define buffer identity,
\path{normalize_maps}, $C/Q$ validity, thresholds, mask application,
empirical-weight separation, and coverage-region thresholding. \\
Configuration and realized state &
\path{include/citlali/core/config/mapmaking_config.h},
\path{include/citlali/core/config/mapmaking_config_validation.h},
\path{include/citlali/core/config/config_error.h},
\path{include/citlali/core/engine/detail/mapmaking_config_impl.h},
\path{include/citlali/core/pipeline/mapmaking_config_policy.h},
\path{include/citlali/core/pipeline/mapmaking_execution_plan.h},
\path{include/citlali/core/pipeline/mapmaking_config_serialization.h}, and
\path{include/citlali/core/pipeline/mapmaking_provenance.h}: define defaults,
validation, one-way adaptation, plan states, serialization, and publication. \\
Callers and reachability &
The direct pointing, Lali, fruit-loop, Beammap population, pipeline, setup,
and mapmaking-pass implementation headers listed in the evidence manifest,
plus \path{include/citlali/core/timestream/rtc/rtcproc.h}: determine ordinary
dispatch, normalization call sites, detector grouping, unique detector-map
indices, outer-scan concurrency, and detector-parallel reachability. \\
Kernel and consumers &
\path{include/citlali/core/timestream/rtc/kernel.h},
\path{include/citlali/core/timestream/timestream.h}, and
\path{include/citlali/core/mapmaking/wiener_filter.h}: determine the
sample-domain template, fruit-loop feedback, and post-map consumer use of
kernel, weight, and coverage.  Filtering mathematics itself was not audited. \\
Product publication and identity &
The direct FITS image helper/name/unit/metadata/write-slot headers, observation
and coadd map-file helpers, and Engine output implementation headers listed in
the evidence manifest: determine array/file identity, WCS attachment, HDU
names, units, formal/empirical labels, coverage mask construction, and
required-output failure behavior. \\
Coadd boundary &
\path{include/citlali/core/pipeline/observation_coadd_accumulation.h},
\path{include/citlali/core/pipeline/raw_observation_outputs.h},
\path{include/citlali/core/pipeline/raw_coadd_outputs.h}, and their direct
execution wrappers: distinguish a JINC observation estimator from the later
weighted coadd without auditing the general coadd contract. \\
\bottomrule
\end{longtable}

\section{Independent reference contract}

The frozen core defines the peak-normalized family
\begin{equation}
  \jinc(z)=\begin{cases}2J_1(z)/z,&z\ne0,\\1,&z=0,\end{cases}
\end{equation}
and requires the effective analytic family, radial coordinate, envelope,
support boundary, pixel-response rule, and array-specific parameter ordering
to be explicit.  Its ideal deposition coefficient is a finite,
dimensionless pixel response
\begin{equation}
  c_{sp}=\frac{1}{A_p}\int_{P_p}g_a(\boldsymbol{x},\boldsymbol{x}_s)
  k_a(\rho_a;\boldsymbol{\theta}_a)\,d^2\boldsymbol{x}.
\end{equation}
The core allows negative and zero coefficients.  It distinguishes geometric
exposure $T_p=\sum_s\Delta t_s$ from coefficient and statistical support.

For mutually independent samples with absolute inverse-variance weights
$q_s$, Equations~\eqref{eq:raw}--\eqref{eq:final} give
\begin{equation}
  \operatorname{Var}_{\mathrm{formal}}(\widehat m_p)=Q_p/C_p^2.
\end{equation}
Valid support requires finite accumulators, positive geometric exposure,
$Q_p>0$, $C_p\ne0$, and any stronger declared conditioning/coverage gate.
The realized response to source template $b_a$ is
\begin{equation}
  R_{a,p}(\boldsymbol{x}_0)=
  \frac{\sum_s q_sc_{sp}b_a(\boldsymbol{x}_s-\boldsymbol{x}_0)}{C_p}.
  \label{eq:response}
\end{equation}
The persisted kernel must be the realized response in signal units, not a bare
coefficient, $C_p$, or $Q_p$.  Formal standardized signal and empirical
signal-to-noise are distinct products.  Stable array identity is one of
\texttt{a1100}, \texttt{a1400}, and \texttt{a2000}, never a positional
fallback.  Requested, effective, observation-resolved, and realized state are
separate one-way records.

\section{Source model and conformity}

\subsection{Analytic family, footprint, and subpixel rule}

\srclabel{} At nonzero physical angular radius $r$, the implementation first
uses $r/l_{d,a}$ and evaluates
\begin{equation}
 k_a(r)=
 \frac{2J_1(2\pi r/(a l_{d,a}))}{2\pi r/(a l_{d,a})}
 \exp\!\left[-\left(\frac{2r}{b l_{d,a}}\right)^c\right]
 \frac{2J_1(3.831706 r/(r_{\max}l_{d,a}))}
      {3.831706 r/(r_{\max}l_{d,a})};
 \label{eq:implkernel}
\end{equation}
at exactly zero it returns one.  Shape-vector positions are $(a,b,c)$.
The fixed numeric scales map keys 0, 1, and 2 to wavelengths 1.1, 1.4, and
2.0 mm divided by 45.  Evidence:
\path{jinc_mm.h:117-160,214-224} and
\path{mapmaking_method_config.h:30-54,86-98}.

\srclabel{} The cache half-width is
$R_a=\lfloor r_{\max}l_{d,a}/\mathrm{pixel\_size\_rad}\rfloor$, and every
cell in the full $(2R_a+1)^2$ square is evaluated.  There is no radial support
predicate; corners extend to approximately $\sqrt{2}\,r_{\max}l_{d,a}$.
The same $r_{\max}$ also places the first zero of the second JINC factor.
Map-edge blocks are cropped without wrap or reflection, but a sample whose
rounded center is outside the map is discarded even when part of its square
footprint could overlap.  Evidence: \path{jinc_mm.h:125-128,162-207,530-553,
1078-1113}.

\srclabel{} For \path{subpixel_n}$=n>1$, $n^2$ point-evaluated phase caches
use offsets $-1/2+(i+1/2)/n$.  The sample coordinate is rounded, the residual
phase is binned, and one cache is selected.  This converges to point sampling
at the realized phase, not to a pixel-area integral.  Evidence:
\path{jinc_mm.h:140-153,185-207,496-515,1034-1061}.

\inflabel{} The implementation is mathematically explicit, but the names
\path{r_max} and \path{subpixel_n} do not by themselves disclose that
the former is both a square-cache scale and a second-JINC zero or that the
latter refines a point-phase approximation.  This is a direct divergence from
the frozen radial/pixel-response contract, not evidence that the mature
parameter values should be optimized.

\subsection{Configuration domain and array identity}

\srclabel{} Typed configuration defaults provide all three stable array keys
and three positional values.  Validation requires finite positive
\path{r_max}, \path{subpixel_n}$\ge1$, and finite shape elements, but it
does not require the first two shape scales to be positive or otherwise
validate the domain of $(a,b,c)$.  Deposition checks finite signal and finite
positive $q_s$ but does not check finite $c_{sp}$.  Evidence:
\path{mapmaking_config.h:40-47},
\path{mapmaking_config_validation.h:32-47},
\path{config_error.h:108-139}, and
\path{jinc_mm.h:488-495,1034-1041,1098-1190}.

\inflabel{} A finite but zero or physically invalid scale can reach
Equation~\eqref{eq:implkernel}; any resulting non-finite coefficient can
contaminate raw $N,C,Q$, kernel, and coverage before the final mask.  Static
source is sufficient to establish the missing admission and coefficient
check.  The exact floating result for every invalid tuple is deliberately not
claimed without execution.

\srclabel{} Request parsing selects shape vectors by stable array name, and a
one-way adapter maps them to fixed numeric array IDs.  Calibration defines
\texttt{0->a1100}, \texttt{1->a1400}, and \texttt{2->a2000}; deposition uses
the detector APT array value, and FITS filenames/PHDU/WCS use the same stable
map.  Missing required names fail through \texttt{at()} rather than neighbor
fallback.  Evidence: \path{engine/calib.h:74-79},
\path{mapmaking_method_config.h:30-54,86-98},
\path{map_observation_files_impl.h:39-53}, and
\path{map_image_output_impl.h:21-50}.

\inflabel{} The current fixed identity mapping conforms for array reordering
and subsetting because selection remains keyed by numeric physical array ID,
not container position.  Runtime evidence for arbitrary malformed APT
identity is outside this package.

\subsection{Normalization, cancellation, and formal weight}

\srclabel{} Accepted samples update $N,C,Q$ exactly as
Equation~\eqref{eq:raw}; the same coefficient sign is used in $N$ and $C$,
while $Q$ uses $c^2$.  Evidence:
\path{jinc_mm.h:559-615,1119-1175}.

\srclabel{} The grid-weight finalizer requires finite
$|C_p|>10^{-8}$ and finite $Q_p>0$, computes
\begin{equation}
 W_p=\frac{C_p^2}{\max(Q_p,10^{-30})},
\end{equation}
then applies a positive-weight mask and, when configured, a support threshold
derived at \path{coverage_cut/10}.  It publishes $N_p/C_p$, $K_p/C_p$,
normalized noise numerators, and the finalized weight on that mask.  It clears
the raw \path{grid_weight} afterward.  Evidence:
\path{map.cpp:76-121,223-301}.

\inflabel{} Negative $C_p$ is valid if its magnitude passes the gate; exact
and sufficiently small cancellation is masked.  Near cancellation above the
fixed gate lowers $C_p^2/Q_p$ as required.  Therefore the signed estimator and
the ordinary cancellation direction conform.

\inflabel{} The two fixed numerical constants carry implicit units.  Under
$d\mapsto\lambda d$ and $q\mapsto q/\lambda^2$, the estimator and formal
weight scale correctly only while the same pixels stay on the same side of
$|C|=10^{-8}$ and $Q=10^{-30}$.  The validity boundary can change solely with
the declared signal unit, and $0<Q<10^{-30}$ does not receive the exact
$C^2/Q$ weight.  Neither constant nor the realized threshold is present in
mapmaking provenance.  Thus final normalization is \emph{conditionally}
conforming, not exact over the admitted numeric domain.

\subsection{Kernel, coverage, masks, and product identity}

\srclabel{} JINC accumulates
\begin{equation}
 K_p=\sum_s q_sc_{sp}h_s,\qquad
 H_p=\sum_s c_{sp}^2/f_s,
\end{equation}
into \path{kernel} and \path{coverage}, respectively.  The finalizer
publishes $K_p/C_p$ and masks $H_p$ but does not otherwise normalize coverage.
Evidence: \path{jinc_mm.h:617-627,1177-1190} and
\path{map.cpp:244-295}.

\srclabel{} Analytic Gaussian and Airy RTC paths describe and create a
unit-source, unit-peak sample-domain template $h_s$; FITS template mode copies
image values and therefore inherits the input image's amplitude convention.
Evidence: \path{timestream/rtc/kernel.h:172-223,226-360,363-410}.
The output writer stores \path{kernel_*} in the signal unit with the same
WCS, plus kernel type and FWHM.  Fruit-loop feedback samples this stored map,
and the Wiener path directly uses it as the filter template/response input.
Evidence: \path{map_image_output_helpers.h:76-95},
\path{fits_image_metadata_keys.h:124-128},
\path{timestream.h:2672-2697,2817-2844}, and
\path{wiener_filter.h:235-255,1324-1349}.

\inflabel{} For the visible analytic unit-source templates, $K/C$ has the
realized-response form in Equation~\eqref{eq:response} and is consumed rather
than reconstructed downstream.  This is a conforming kernel equation.  The
FITS-template amplitude convention and complete source identity are not
carried by the kernel HDU or mapmaking provenance, so product provenance is
only partial.  A unit-source runtime response comparison remains unverified.

\srclabel{} \path{coverage_*} is labeled seconds and ``effective integration
time,'' while its actual estimator is $H_p=\sum c^2/f_s$.  It is not geometric
exposure, contributor time, $C$, or $Q$; analytic coefficient zeros count as
zero coverage.  The Wiener edge guard may use this plane to refine support.
Evidence: \path{fits_image_units_kernels.h:33-35,116-121} and
\path{wiener_filter.h:1074-1105}.

\srclabel{} The writer does not derive \path{coverage_bool_*} from
\path{coverage}.  It obtains a threshold from the current working weight
and returns one wherever \texttt{weight < threshold} is false.  If the
threshold is zero, a zero-weight pixel is therefore labeled one; a non-finite
weight also compares false and is labeled one for any threshold.  The default
\path{coverage_cut} is zero.  Evidence:
\path{map_image_output_helpers.h:98-123},
\path{fits_image_products.h:65-100},
\path{map.cpp:539-551}, \path{utils.h:526-595}, and
\path{mapmaking_config.h:208-228}.

\inflabel{} This directly falsifies the published-mask contract: positive
weight and finite signal-like state are not required for a pixel labeled
usable.  It is not merely an unverified edge case.  Separately, the coverage
HDU has a dimensionally valid second unit but an incompletely declared
coefficient-squared estimator; consumers cannot infer geometric exposure from
it.

\srclabel{} Signal, working weight, formal weight, noise variance, kernel,
coverage, coverage mask, formal standardized signal, and empirical pixel S/N
have distinct HDU names and metadata.  When empirical noise products exist,
the code saves the pre-calibration working weight as \path{weight_formal_*},
derives a map-level empirical scale from jackknife variance, optionally
replaces \path{weight_*}, and labels pixel S/N empirical.  Without empirical
products, signal times square-root weight is labeled
\path{formal_standardized_signal_*}, explicitly not significance.
Evidence: \path{map.cpp:826-923},
\path{map_image_output_helpers.h:13-73,125-155},
\path{fits_image_hdu_names_wcs.h:5-53}, and
\path{fits_image_units_kernels.h:19-27,124-137}.

\inflabel{} The formal-versus-empirical standardized-product naming conforms
for the traced observation product.  Empirical-noise estimator validity is
outside \packageid{} and is not audited here.

\subsection{Observation/coadd boundary}

\srclabel{} Observation normalization occurs before raw observation noise
products and output.  Optional empirical calibration can replace the working
observation weight.  Coaddition later accumulates
\begin{equation}
 \sum_o W_{o,p}\widehat m_{o,p},\qquad \sum_o W_{o,p},
\end{equation}
and likewise weight-averages the normalized kernel; it sums coverage.  The
coadd finalizer then uses its non-grid-weight branch.  Evidence:
\path{raw_observation_outputs.h:13-49},
\path{observation_output_execution.h:34-57},
\path{observation_coadd_accumulation.h:19-93}, and
\path{raw_coadd_outputs.h:21-45}.

\inflabel{} The per-observation JINC product is the direct $N/C$ estimator.
The coadd is a later inverse-working-weight mean of already normalized
observations, and the chosen $W_o$ may be formal or empirically rescaled.  It
is not generally the single raw-sample ratio $(\sum_oN_o)/(\sum_oC_o)$.
This report records the identity boundary but does not reopen the excluded
general coadd audit or reinterpret \path{SCI-MAP-001} evidence.

\subsection{Requested, effective, resolved, and realized provenance}

\srclabel{} \path{mapmaking_provenance.yaml} serializes the complete requested
and effective mapmaking configurations, including named JINC vectors;
requested/effective grouping and unit resolution; observation index, obsnum,
map count, effective pixel size, write count and completion; coadd
cardinality; and coarse realized completion/mapmaking booleans.  It is written
atomically after lifecycle completion.  Evidence:
\path{mapmaking_config_serialization.h:15-149},
\path{mapmaking_provenance.h:13-49},
\path{mapmaking_execution_plan.h:18-179}, and
\path{reduction_execution.h:260-270}.

\inflabel{} This is a sound one-way state skeleton, but it omits the analytic
JINC convention, $(a,b,c)$ semantics, square support, realized pixel cutoff,
present array-to-parameter selections, coefficient/cancellation diagnostics,
normalization constants and support threshold, execution policy, kernel
source/amplitude convention, coverage estimator, product paths/digests, and
formal/empirical association.  It therefore cannot prove which realized JINC
contract produced a downstream map.

\subsection{Sequential and parallel reachability}

\srclabel{} Ordinary typed dispatch calls \path{populate_maps_jinc}.  That
method accumulates into thread-local scratch and merges touched blocks under a
mutex.  Non-detector Beammap grouping can run scans under \texttt{grppi} and
calls this same method.  Evidence: \path{mapmaking_dispatch.h:9-27},
\path{jinc_mm.h:333-425,729-775}, and
\path{beammap_map_population_impl.h:43-54}.

\srclabel{} Detector-grouped Beammap directly calls
\path{populate_maps_jinc_parallel}, which parallelizes over detectors and
writes map blocks without a numerical lock.  However, detector grouping
constructs \path{indices=LinSpaced(0,n_dets-1)} and then a one-to-one
detector-to-map mapping; therefore each reachable detector worker writes a
different map plane.  Evidence:
\path{beammap_map_population_impl.h:17-35},
\path{rtcproc.h:791-846}, and
\path{jinc_mm.h:964-1005,1115-1190}.

\inflabel{} The preliminary phase-two ``overlapping write'' concern is not a
current reachable defect under that caller invariant.  The implementation
would be unsafe for a future concurrent caller that maps two workers to the
same plane, but no such direct caller was found in the bounded trace.
Observation/scans can still merge in runtime-dependent order, so bitwise or
tolerance-bounded sequential/parallel equivalence is not established
statically.

\runlabel{} No exact-SHA reduction demonstrates which policy executed, stable
contributor cardinality, cancellation-sensitive agreement, output mask
agreement, or product-level sequential/parallel acceptance.

\section{Equation/source conformance table}

\small
\begin{longtable}{@{}p{0.18\textwidth}p{0.27\textwidth}p{0.27\textwidth}p{0.20\textwidth}@{}}
\toprule
Contract item & Frozen equation/identity & Governing source & Disposition \\
\midrule
\endhead
Zero-limit JINC & $\jinc(0)=1$ & Exact \texttt{r==0} return one; nonzero path
uses two JINCs and envelope. & Conforms at the removable limit; analytic family
must be documented as the product in Equation~\eqref{eq:implkernel}. \\
Footprint & Declared radial $g_{sp}$ & Full square cache; no radial cutoff;
$r_{\max}$ also second-JINC zero. & Divergent from frozen radial support (F002). \\
Pixel response & Pixel-area average or declared convergent quadrature & One
phase-quantized point cache & Divergent target; scientifically undecided
whether contract or implementation should change (F003). \\
Numerator & $N=\sum qcd$ & \texttt{signal += q*c*d} & Conforms. \\
Signed denominator & $C=\sum qc$ & \path{grid_weight += q*c} & Conforms. \\
Variance term & $Q=\sum qc^2$ & \texttt{weight += q*c*c} & Conforms. \\
Signal finalization & $\widehat m=N/C$ on valid support & \texttt{signal/C}
with $|C|>10^{-8}$ and mask & Conforms algebraically; fixed absolute gate is
unit-dependent (F004). \\
Formal weight & $C^2/Q$ & $C^2/\max(Q,10^{-30})$ after $Q>0$ & Conforms for
$Q\ge10^{-30}$; conditional divergence below the floor (F004). \\
Signed cancellation & Exact $C=0$ invalid; near cancellation low weight &
Negative $C$ accepted; $|C|\le10^{-8}$ masked; otherwise $C^2/Q$ & Direction
conforms; absolute scale and provenance do not (F004). \\
Coefficient validity & Finite $c$ or fail before deposition & Shape elements
finite only; no physical-domain or coefficient-finite check & Nonconformant
(F005). \\
Kernel response & $R=\sum qch/C$ in signal units & $K=\sum qch$, then $K/C$;
analytic $h$ is unit source; same map consumed & Equation conforms; FITS
amplitude/provenance and runtime injection remain open (F007/F008). \\
Coverage & Declared geometric exposure $T=\sum\Delta t$ &
$H=\sum c^2/f_s$, seconds, then final mask & Different effective-exposure
estimator and incomplete metadata (F006). \\
Published support & Explicit finite/positive validity authority &
\path{coverage_bool=(weight<threshold)?0:1} & Zero/non-finite weights can be
labeled valid; nonconformant (F001). \\
Formal/empirical identity & Separate formal weight, empirical weight/SNR,
formal standardized plane & Separate HDUs and estimator metadata; formal
standardized signal only without empirical products & Conforms for traced
observation publication. \\
Array identity & Stable name selects shape and product & Fixed ID/name map,
name-keyed config/provenance, array-keyed filename/WCS & Conforms for valid
current calibration identity. \\
Four-stage state & Lossless requested/effective/resolved/realized JINC and
product join & Full requested/effective config; coarse observation and realized
cardinality only & Partial/nonconformant provenance (F007). \\
Parallel equivalence & Same contributors/equations and accepted numerical
policy & Same equations; detector-parallel maps are disjoint; outer-scan merge
order may differ & Static race concern discharged; numerical equivalence
unverified (F008). \\
\bottomrule
\end{longtable}
\normalsize

\section{Findings}

\begin{longtable}{@{}p{0.14\textwidth}p{0.09\textwidth}p{0.69\textwidth}@{}}
\toprule
ID & Severity & Finding and closure condition \\
\midrule
\endhead
\texttt{F001} & P0 & \textbf{Published validity mask admits invalid support.}
\srclabel{} \path{coverage_bool} uses a strict less-than comparison against
the weight threshold.  With the default zero coverage cut, zero weight is
labeled one; non-finite weight is also labeled one because the comparison is
false.  Closure requires an approved finite/positive mask predicate, exact
HDU semantics, zero/NaN/Inf and threshold-boundary tests, exact-SHA product
evidence, and recipient re-audit. \\
\texttt{F002} & P1 & \textbf{\texttt{r\_max} is not the frozen radial support
cut.}  It is both a second-JINC zero scale and half-width of a fully populated
square.  Closure requires the scientific owner to approve either the realized
square/dual-use convention or a radial support contract, serialize it, test
below/equal/above and corner/edge cases, and re-audit; no parameter optimization
is implied. \\
\texttt{F003} & P1 & \textbf{Subpixel refinement targets point-phase sampling,
not pixel response.}  Closure requires an owner decision defining the target
response, phase/rounding convention and convergence criterion, followed by
bounded tests and exact-successor-SHA re-audit. \\
\texttt{F004} & P1 & \textbf{Final normalization is unit-scale dependent at
fixed absolute gates.}  $|C|>10^{-8}$ and $Q$ floor $10^{-30}$ are not expressed
in declared units or provenance and can change support or weight under a unit
rescaling.  Closure requires an approved conditioning policy, typed units,
serialization, scaling/cancellation/tiny-$Q$ tests, and re-audit. \\
\texttt{F005} & P1 & \textbf{Parameter admission and coefficient finiteness are
incomplete.}  Finite shape values with zero/invalid scales are admitted, and
non-finite evaluated coefficients are not rejected before accumulation.
Closure requires an approved $(a,b,c)$ domain and fail-closed coefficient
policy plus malformed-domain tests; it does not require running invalid
reductions on Unity. \\
\texttt{F006} & P1 & \textbf{Coverage identity is coefficient-squared exposure,
not geometric exposure.}  The seconds-valued $\sum c^2/f_s$ plane loses
analytic-zero geometric exposure and is consumed as an edge-support proxy,
while metadata states only ``effective integration time.''  Closure requires
an owner-approved estimator name/accounting rule, explicit distinction from
geometric exposure and masks, product metadata/provenance, consumer admission,
and exact-SHA verification. \\
\texttt{F007} & P1 & \textbf{Realized JINC and product provenance is
incomplete.}  Requested/effective config exists, but realized analytic,
support, normalization, array-selection, execution, kernel-source, coverage,
cancellation, mask, and product-link identities do not.  Closure requires a
versioned one-way record and exact serialization/product-join tests. \\
\texttt{F008} & P1 & \textbf{No exact-governing-SHA JINC runtime evidence is
available.}  The incoming evidence is naive-only.  Static analysis resolves
the current detector-parallel writes as map-disjoint but cannot establish
response, product values, execution-policy realization, or numerical
sequential/parallel acceptance.  Closure requires coordinator-approved return
of the bounded prerequisites in Section~\ref{sec:protocol}; no evidence is
requested by this report. \\
\bottomrule
\end{longtable}

\section{Decisions and dependencies}

\begin{enumerate}[leftmargin=*]
  \item \textbf{SCI-MAP-002-D001 --- amend.}  Adopt
        Equations~\eqref{eq:raw}--\eqref{eq:final} as the proposed JINC
        observation estimator, while recording the current fixed numerical
        validity policy as nonconforming until scientifically resolved.
  \item \textbf{SCI-MAP-002-D002 --- restrict.}  Do not use
        \path{coverage_bool_*} as an authoritative validity mask for new
        scientific validation or broadened production claims until F001 is
        closed.  This audit does not alter existing runs.
  \item \textbf{SCI-MAP-002-D003 --- owner decision required.}  The square
        footprint, dual-use \path{r_max}, point-phase subpixel response, and
        coefficient-squared coverage need explicit scientific acceptance or a
        separately authorized successor implementation.  No parameter study
        is authorized; any such study routes through \texttt{FRAMEWORK-NUM-001}.
  \item \textbf{SCI-MAP-002-D004 --- retain positive controls.}  Preserve the
        distinct $N,C,Q$ accumulators, signed coefficients, $K/C$ response
        equation, stable array-name mapping, and formal/empirical product
        separation in any future work.
  \item \textbf{SCI-MAP-002-D005 --- evidence remains proposed.}  The protocol
        below is a prerequisite specification only.  It is not an external
        evidence request and does not authorize Codex to access Unity.
\end{enumerate}

The package retains these dependencies:

\begin{itemize}[leftmargin=*]
  \item \textbf{SCI-MAP-001/general coadd:} this report distinguishes the
        observation-level JINC estimator from the later formal/empirical
        weighted coadd.  It neither closes nor reopens SCI-MAP-001 findings.
  \item \textbf{SCI-NOI-001/002 and SCI-FLT-001/002:} downstream conclusions
        remain conditional on a future accepted JINC product contract.  The
        visible consumers use kernel/weight/coverage directly, so F001, F006,
        and F007 require recipient admission/re-audit before closure.
  \item \textbf{Scientific JINC owner:} must decide the footprint, subpixel,
        conditioning, coverage, and parameter-domain contracts before a repair
        or parameter campaign can be scoped.
\end{itemize}

\section{Proposed exact-SHA Unity evidence protocol}
\label{sec:protocol}

This section is a \emph{proposal}, not a request.  Nothing here was run or
sent to Unity.  A coordinator must separately approve the inputs and human
execution.

\subsection{Prerequisites returned with every case}

\begin{enumerate}[leftmargin=*]
  \item A clean checkout whose application identity is exactly
        \texttt{\appsha}; return \texttt{git status}, exact SHA, compiler,
        build mode, dependency lock identity, host/CPU, thread environment,
        and executable digest.
  \item A coordinator-approved, checksum-pinned input observation or existing
        simulation with explicit reduction type, present arrays, detector
        identity, sample rate, flags, and calibration identity.  A true
        unit-source response claim additionally requires a checksum-pinned
        unit-amplitude input whose source template is independently declared;
        if none exists, F19-style response remains open rather than launching
        a campaign.
  \item The exact authored YAML and its checksum, complete resolved config,
        \path{mapmaking_provenance.yaml}, stdout/stderr, runtime thread/policy
        diagnostics, product index, and every returned FITS file/digest.
  \item For each FITS product, return HDU names, shapes, numeric type, WCS,
        UNIT/BUNIT/TYPE/DESCRIP/WTTHRESH/FWHM keys, and finite/NaN/Inf/zero/
        positive pixel counts.  Preserve both working and formal weights when
        empirical noise products are enabled.
  \item No code patch, helper, validator, parameter optimization, or repair in
        the evidence branch.  An unavailable prerequisite yields ``not
        established,'' not substitution by a different SHA or naive run.
\end{enumerate}

\subsection{Bounded case set}

\begin{longtable}{@{}p{0.10\textwidth}p{0.29\textwidth}p{0.53\textwidth}@{}}
\toprule
Case & Controlled difference & Required comparison \\
\midrule
\endhead
U01 & One ordinary array-grouped JINC reduction with raw kernel enabled and
empirical replacement disabled. & Confirm method=JINC and present-array
parameter selection in resolved config; return signal, formal weight, kernel,
coverage and mask products; verify product/header joins and finite-positive
mask implications.  This is product evidence, not parameter validation. \\
U02 & Same input/config as U01, differing only in empirical-weight replacement
enabled with fixed noise settings. & Confirm signal/kernel/coverage identity
unchanged within the approved policy; working weight and empirical S/N change
as declared; formal weight remains separately identical; record any coadd
weight identity separately. \\
U03 & Detector-grouped Beammap JINC, forced sequential execution policy. &
Return policy realization, active detector-to-map mapping, normalized products,
support/cancellation diagnostics, logs, and digests. \\
U04 & Bit-for-bit same as U03 except the authorized concurrent policy and
thread count. & Compare contributor/map cardinality, valid-mask set, all
finite-status counts, and signal/formal-weight/kernel/coverage arrays using a
coordinator-approved absolute/relative policy with special reporting for
small $|C|$ support.  Do not assume bitwise equality. \\
U05 & Repeat U04 once under the same policy and thread environment. & Establish
or falsify run-to-run determinism at the named policy; report array-level
max/quantile differences and changed support, never only file hashes. \\
U06 & Only if the checksum-pinned independent unit-source input in prerequisite
2 already exists: run the corresponding JINC case with the exact template and
flags. & Compare output signal response to the recorded kernel on their common
finite valid support, including peak, integral, edge asymmetry, WCS, units and
array identity.  Otherwise mark ``not available'' and do not create a new
simulation campaign under this package. \\
\bottomrule
\end{longtable}

The evidence acceptance policy must be approved before execution and must
name tolerances, valid-pixel matching, non-finite handling, and cancellation-
sensitive reporting.  The bundle may inform re-audit only after its exact
identity and completeness are checked; it cannot itself amend the contract.

\section{Limits of this audit}

This report does not assess naive or maximum-likelihood mapmaker internals,
the general coadd repair, RTC/PTC scientific correctness, empirical noise
validity, Wiener/convolution mathematics, calibration, astrometry, flagging,
source fitting, Beammap inference, or JINC parameter optimality.  Direct
consumer lines were opened only to determine product identity and reachability.
No test result or Unity behavior is implied by static source.  No source
change follows automatically from a finding.  A future repair, validation
run, evidence request, downstream audit, or numerical study requires separate
coordinator authorization.

\end{document}
